\section{Introduction}
A signed release manifest lets an auditor authenticate a developer's description of a model. A harder question arises when a third party distributes a modified checkpoint: can an auditor who holds the original manifest still connect the weights to that release? Exact hashing answers whether two files are identical. A provenance watermark addresses a different question, because fine-tuning or quantization can change the checkpoint while preserving some of its learned behavior.

We investigate this question on a CIFAR-10 convolutional classifier. The developer signs a record of the release and embeds a fingerprint of that record in the weights. The signature authenticates the record, while a keyed measurement of one convolutional layer supplies evidence from the checkpoint. These steps have different failure modes: a valid signature can accompany an unrelated model, and a damaged watermark can carry statistical evidence even when decoding fails.

Our experimental question is whether reducing the classification cost of embedding preserves resistance to removal. We vary the embedding loss while holding the signed payload, carrier layer, and error-correcting code fixed. We compare accuracy across matched training seeds, then evaluate three watermark configurations under the same checkpoint modifications. The distinction between exact decoding, residual evidence, and successful removal is central to the comparison: the configuration with higher clean accuracy can also be easier to overwrite.
